% Generated draft: requires mathematical and source review.
\documentclass[11pt]{article}
\usepackage[margin=25mm]{geometry}
\usepackage{fontspec}
\setmainfont{Latin Modern Roman}
\usepackage{amsmath,amssymb,mathtools,unicode-math}
\setmathfont{Latin Modern Math}
\usepackage{enumitem,needspace,xurl,hyperref}
\setlength{\parindent}{0pt}
\setlength{\parskip}{5pt}
\setlength{\emergencystretch}{3em}
\newcommand{\N}{\mathbb{N}}
\newcommand{\Z}{\mathbb{Z}}
\newcommand{\Q}{\mathbb{Q}}
\newcommand{\R}{\mathbb{R}}
\newcommand{\C}{\mathbb{C}}
\newcommand{\F}{\mathbb{F}}
\newcommand{\A}{\mathbb{A}}
\newcommand{\PP}{\mathbb{P}}
\newcommand{\E}{\mathbb{E}}
\newcommand{\eps}{\varepsilon}
\newcommand{\dd}{\,\mathrm{d}}
\begin{document}
\section*{UnsolvedMath ID: 4800027}
\textbf{Source problem code:} catalogue entry
\subsection{Definitions and mathematical statement}
Let $p_n$ be the $n$th prime, in increasing order. Fix a real number $c>0$ with $c\notin\Z$, and put $a(n)=\lfloor p_n^c\rfloor$. The sequence is indexed by primes; repeated values of $a(n)$ are retained.

The source asks that this sequence be good for multiple recurrence and convergence of powers. Under the usual single-transformation norm-convergence interpretation, for every integer $\ell\geq1$, invertible probability-preserving system $(X,\mathcal B,\mu,T)$, and bounded measurable functions $f_1,\ldots,f_\ell$, the averages
\[
\frac1N\sum_{n=1}^N\prod_{j=1}^{\ell}f_j\bigl(T^{j\lfloor p_n^c\rfloor}x\bigr)
\]
should converge in $L^2(\mu)$ as $N\to\infty$.

The corresponding recurrence statement concerns every $\ell\geq1$ and every $A\in\mathcal B$ with $\mu(A)>0$, and asks for return patterns
\[
\mu\left(A\cap T^{-a(n)}A\cap\cdots\cap T^{-\ell a(n)}A\right)>0.
\]
One common quantitative convention requires
\[
\liminf_{N\to\infty}\frac1N\sum_{n=1}^N
\mu\left(A\cap T^{-a(n)}A\cap\cdots\cap T^{-\ell a(n)}A\right)>0;
\]
a weaker convention requires only existence of an index $n$. The record does not reproduce the survey's definition of ``good'', so the precise recurrence quantifier must be checked rather than silently chosen. Likewise, convergence is displayed in its usual norm sense, not as an asserted pointwise result. Powers here are multiples of one transformation, not independent commuting transformations; that convention also requires source checking.
\subsection{Short English statement}
For every positive noninteger $c$, establish multiple recurrence and norm convergence along the integer parts of the $c$th powers of primes, with the exact ``good'' conventions of the cited survey.
\subsection{Sources}
\begin{itemize}
\item Local UnsolvedMath registry, permanent ID 4800027, exact source code catalogue entry, ``Multiple ergodic averages---Problem 27''; proposed by Nikos Frantzikinakis, 2016. Registry status solved, classification SOLVED-IN-LITERATURE.
\item Nikos Frantzikinakis, \emph{Some open problems on multiple ergodic averages} (2016), Problem 27, \url{https://arxiv.org/abs/1103.3808}. Supplied source-TeX extraction accessed July 29, 2026.
\item The research summary attributes a solution for fractional powers of primes to Koutsogiannis--Tsinas, ``to appear J. Modern Dynamics'', with support cited only as ``N.F. Sigma 2022''. It supplies no titles, theorem statements, or full bibliographic details. It also says recurrence for $0<c<1$ was previously known by a finite-miss-range argument. These are supplied claims, not independently checked.
\item Supplied provenance says the author's November 2025 progress page does not record a resolution of this numbered problem and flags status NEEDS\_REVIEW; no page URL is supplied. It flags rights NEEDS\_REVIEW because public access did not establish redistribution permission.
\end{itemize}
\end{document}
